\section{Security and approximation boundaries}\label{sec:security}
\subsection{Threat model and two-stage binding}
The prover may choose false statistics, assignments, costs or authentication paths. The verifier has an authentic original root, approved public model and verification-key hashes. We assume binding commitments, collision-resistant context hashing, and knowledge soundness and zero knowledge of the instantiated proof system. Groth16 setup is trusted; a compromised setup falls outside these assumptions. The proof establishes computations on the owner's committed catalog, not the truth or provenance of the owner's measurements.

\begin{theorem}[Binding of reused integer audits]\label{thm:binding}
Suppose the verifier accepts complete Stage~1 coverage for one context and accepts a Stage~2 proof under the pinned keys. It also checks that the public new-center statistics are the approved integer derivations of the intended centers. Except for a commitment/hash collision or a proof-system soundness failure, the accepted cost and margin equal Eq.~\eqref{eq:reuse-outputs} on the deterministic statistics of the committed original rows at the public requested indices.
\end{theorem}
\begin{proof}
Extract satisfying witnesses from the accepted proofs. Complete Stage~1 coverage binds every original index to its derived statistic row under the common context. Distinct openings under the same authenticated roots would contradict binding. Thus every row opened in Stage~2 is the same derived row established in Stage~1. The range checks give the specified integer semantics; the score comparisons enforce first-index minimizers and the second-smallest distance. The output constraints then give the claimed sums. The verifier's center-statistic check identifies those public inputs with the intended centers. No approximation claim is needed for this argument.
\end{proof}

Knowledge soundness applies to the proof transcript; the union of soundness errors across a finite collection of proofs and commitment/hash failures bounds failure of the combined audit. This does not certify closeness to floating-point Full. That is determined by basis coverage, the omitted cross term and quantization. In particular, a malicious choice of an approved learned basis can give a faithfully proved but poor approximation. Orthonormality, center-span coverage, basis training policy and approval remain external checks.

\subsection{Public outputs and verification boundary}
Original rows, derived heads/energies, paths, salts and assignments remain private witnesses. The public statement exposes roots, model metadata, subset indices, center statistics, aggregate cost and aggregate margin. Diagnostic assignment arrays are saved by the experimental prover for fidelity analysis; they are not required public inputs. Repeated costs and margins can reveal information about the private data, and zero knowledge does not provide differential privacy or an unrestricted-query leakage bound. Owner/model approval and a query policy are application responsibilities.

Every generated benchmark proof is verified by the proving CLI. The saved original fixed-center bundles also support public-only re-verification from pinned keys and public statements. Reuse requires the verifier to supply the intended new centers and recompute their public statistics; acceptance of a prover-supplied cost alone is insufficient. The artifact validation records which of these checks were actually performed for each cohort.

\subsection{Commit-then-sample JL option}\label{sec:theory}
A data-oblivious alternative fixes data and the queried parameter set before fresh projection randomness. Register the original root $C_X$ and a digest $C_\theta$ of the public task, then obtain unpredictable challenge $\rho$ and derive
\begin{equation}
s=H(\texttt{CommitJL-v1}\|C_X\|C_\theta\|\rho).
\end{equation}
The verifier expands a canonical sparse signed map $A$, hashes its canonical encoding, and requires the proved computation to authenticate that same map. Each column has $q$ distinct signed buckets; the real normalized map is $A/\sqrt q$. Canonical sampling and its distribution assumptions are specified in Appendix~\ref{app:jl}. The learned head does not acquire a geometric guarantee merely because a JL tail is attached.

For a point/center set fixed independently of $A$, let $\mathcal E_\epsilon$ be the event that every relevant squared distance is preserved within $1\pm\epsilon$. Conditional on this event, the normalized projected nearest-center cost obeys
\begin{equation}
(1-\epsilon)\Phi_I(\mathcal C)\le
\Phi_{AI}(A\mathcal C)/q\le(1+\epsilon)\Phi_I(\mathcal C).
\end{equation}
Take the minimum of the pointwise bounds, then sum over $I$. A unique nearest center remains nearest if its squared-distance gap satisfies $d_2^2/d_1^2>(1+\epsilon)/(1-\epsilon)$, handling $d_1=0$ separately. Any failure probability requires hypotheses for the actual sparse distribution, dimension and admissible vectors. In particular, the tested one-bucket maps do not receive a dense-JL concentration bound by assumption.

The challenge order prevents choosing data after seeing the random map. It does not address unlimited aborts or recommitments: if a uniform single-attempt geometric failure bound is $\delta$, a policy allowing at most $Q$ eligible attempts admits the union bound $Q\delta$, plus the cryptographic/generator errors. Later centers chosen adaptively after observing the map require a separate argument; the fixed-set JL event does not cover them automatically. Our deployment experiments use locally generated challenges, not a deployed public beacon service.
